\subsection{A completed old tuple with an obstructed rational puncture}

\begin{statement}{A.90}{Proved}
{An ordinary completed tuple that cannot acquire one more coordinate}
Put
\[
 c_1=0,\qquad c_2=1,\qquad c_3=-\frac{9563}{59024},
 \qquad R=\frac1{17},\qquad c=-\frac{11}{68},
 \qquad a_i=R-c_i .
 \tag{A.90.1}
\]
There exist three classes
$A_i\in\widehat{\Q^\times}=\varprojlim_m\Q^\times/\Q^{\times m}$
and, at every finite prime $p$, an ordinary coordinate
$t_p\in\Q_p\setminus\{c_1,c_2,c_3,c\}$ such that
\[
 A_i|_{\Q_p}=\kappa(t_p-c_i)\quad(1\le i\le3),\qquad
 A_i\bmod4=[a_i],\qquad t_2=\frac{16}{17}.
 \tag{A.90.2}
\]
Every $A_i$ has infinitely many nonzero prime valuations.
At the real place the three signs are those of $R-c_i$,
and the ordinary coordinate can be taken to be $R$.

There is no class $A_c\in\widehat{\Q^\times}$ satisfying
$A_c|_{\Q_p}=\kappa(t_p-c)$ at every finite prime.
Already its reduction modulo two is impossible.
The resulting ordinary completed pair
$(A_1,-A_2)$ has local values $(t_p,1-t_p)$;
the third completed coordinate is $A_3$.
Thus the failure of this particular new-coordinate
initialization occurs for actual completed old data,
with the fixed finite local obstruction of A.69.
\end{statement}
\begin{proof}
We construct covers with mixed exponents and exact maps
between them.  The exponents will be cofinal under
divisibility by all positive integers.
The local coordinate $T$ is retained permanently once
a prime enters a finite controlled set.

\emph{The common abelian field for mixed exponents.}
We first prove the version of A.80 needed in this
construction.  Let $4\mid n_j$ for $1\le j\le k$ and put
\[
 \mathcal E=\Q[Z_1^{\pm1},\ldots,Z_k^{\pm1}]
                         /(Z_j^{n_j}-D_j),\qquad
 N=\operatorname{lcm}_j n_j,\qquad
 K=\Q(\sqrt{D_1},\ldots,\sqrt{D_k}).
 \tag{A.90.3}
\]
Assume that $-1$ is outside the squareclass span of
the $D_j$.  We prove that the largest abelian field
contained in every field factor of $\mathcal E$ is
exactly $K$.

Write $D_j=(-1)^{\delta_j}\prod_P P^{w_{jP}}$, using
all positive primes with nonzero valuations in the
finite list of radicands.  Set
\[
\begin{aligned}
 F&=\Q(\mu_{2N},\sqrt P:P),\qquad
 L=F(\alpha_P:P),\qquad \alpha_P=P^{1/N}>0,\\
 G&=\operatorname{Gal}(L/\Q),\qquad
 \Gamma=\operatorname{Gal}(F/\Q).
\end{aligned}
 \tag{A.90.4}
\]
The degree and commutator proof (A.57.6)--(A.57.8),
applied to $N$, gives $L^{\rm ab}=F$ and the following
actual parametrization of the restriction map
$G\twoheadrightarrow\Gamma$.
For $\sigma\in\Gamma$ let $u_\sigma$ be its exponent
on $\mu_{2N}$ and let
$\epsilon_P(\sigma)\in\Z/2$ describe its action on
$\sqrt P$.  Its lifts are exactly the maps
\[
 g(\alpha_P)=\zeta_N^{k_P}\alpha_P,\qquad
 k_P\in\Z/N,\qquad
 k_P\equiv\epsilon_P(\sigma)\pmod2 .
 \tag{A.90.5}
\]
In particular all choices with the stated parities
are available simultaneously.
The root tuples of $\mathcal E$ in $L$ are
\[
 z_j=\zeta_{2n_j}^{\delta_j+2e_j}
             \prod_P\alpha_P^{\,Nw_{jP}/n_j},
                       \qquad e_j\in\Z/n_j.
 \tag{A.90.6}
\]
The lift in (A.90.5) fixes this actual tuple exactly when
\[
 (u_\sigma-1)(\delta_j+2e_j)
                 +2\sum_Pw_{jP}k_P=0\pmod{2n_j}
                                  \quad\text{for every }j.
 \tag{A.90.7}
\]
These formulas include the odd parts of all the
possibly different exponents.

Every tuple stabilizer fixes $K$, since
$z_j^{n_j/2}$ is a square root of $D_j$.
Conversely take $\sigma\in\operatorname{Gal}(F/K)$
with $u_\sigma\equiv3\pmod4$ and write
\[
 d_j=\gcd(n_j,u_\sigma-1)=2d_{j,o},\qquad
 D_o=\operatorname{lcm}_j d_{j,o}.
 \tag{A.90.8}
\]
Here $D_o$ is odd and divides $N$.
Choose each $k_P$ with its prescribed parity and with
$k_P\equiv0\pmod{D_o}$; the Chinese remainder theorem
gives these elements of $\Z/N$.
For each $j$ the integer
\[
       \frac{u_\sigma-1}{2}\delta_j+\sum_Pw_{jP}k_P
 \tag{A.90.9}
\]
is divisible by $d_{j,o}$, and it is even because
$\sigma$ fixes $\sqrt{D_j}$.
It is therefore divisible by $d_j$.
The linear congruence (A.90.7), divided by two, has a
solution $e_j$ for each $j$.
Thus $\sigma$ is the image of an actual tuple stabilizer.

The hypothesis on squareclasses says $i\notin K$.
Choose $\tau\in\operatorname{Gal}(F/K)$ acting nontrivially
on $i$; then $u_\tau\equiv3\pmod4$.
For an arbitrary $\sigma\in\operatorname{Gal}(F/K)$,
either $\sigma$ has this same congruence, or both
$\sigma\tau$ and $\tau$ do.
It follows that the images of all tuple stabilizers
generate exactly $\operatorname{Gal}(F/K)$.
Since every abelian subfield of $L$ lies in $F$,
Galois correspondence proves the asserted common
abelian field:
\[
 F^{\langle\operatorname{im}(G_\omega\to\Gamma):
                                  \omega\text{ a root tuple}\rangle}=K .
 \tag{A.90.10}
\]
This proves the mixed-exponent assertion for all
$4\mid n_j$, including exponents with odd factors.

\emph{The fixed boundary fields and their initial evaluations.}
Start with the curve $C_0$ of A.69,
$T-c_i=a_iX_i^4$.  For each boundary $T=c_i$ define
\[
 B_i=\Q\left(\sqrt{(c_i-c_j)/a_j}:j\ne i\right).
 \tag{A.90.11}
\]
The two radicands in these three fields are
\[
 \begin{array}{c|cc}
 i&\multicolumn{2}{c}{(c_i-c_j)/a_j,\quad j\ne i}\\ \hline
 1&17/16&9563/13035\\
 2&17&68587/13035\\
 3&-9563/3472&68587/55552 .
 \end{array}
 \tag{A.90.12}
\]
The first two rows consist of positive numbers.
In the last row, the first radicand has an odd
valuation at $73$ and the second at $107$, while
the opposite radicand has valuation zero at
the indicated prime.
Consequently no row has $-1$ in its squareclass span.
The mixed-exponent result will therefore apply to
every boundary field in the construction.

The initial local points are the rational point
$P_0=(R,1,1,1)$, except at two, where we use
\[
 P_2=(16/17,2,1/2,r),\qquad
          r^4=4341/869,\quad r\in1+4\Z_2.
 \tag{A.90.13}
\]
The exact fourth-root construction is (A.88.14).
For every character $\chi$ of every $B_i$, these
points have the stronger coordinate-root property
\[
              (X_i(P_v),\chi)=0
                            \quad\text{in }\operatorname{Br}(\Q_v).
 \tag{A.90.14}
\]
At $P_0$ the argument is one.
At two, all four radicands in the first two rows of
(A.90.12) are squares, since their valuations are
even and their normalized units are one modulo eight.
The last row has even valuations and odd-unit
squareclass representatives; every odd unit is a
norm from $\Q_2(\sqrt r)$ when $r\in1+4\Z_2$.
Explicitly the only nontrivial squareclass of such
an $r$ is five, and
$N(a+(1+\sqrt5)/2)=a^2+a-1$ takes every odd-unit
value by simple-root Hensel lifting.
This proves (A.90.14), including its individual
vanishing at each place.
At the real point all root coordinates are positive.

\emph{A fixed finite set supporting the obstructing evaluation.}
Choose a finite set $S_*$ of primes containing two,
the coefficient and puncture-denominator primes in
(A.90.1), all nonzero puncture-difference primes
among $c_1,c_2,c_3,c$, the ramified primes of all
$B_i$, and the primes needed in the outside-set
calculation (A.85.4)--(A.85.6) for $C_0$ and the
new puncture $c$.
Include also the ramification of the constant
quaternion $(R-c,2)$.
Prescribe $P_2$ at two and $P_0$ at every other
prime in $S_*$ and at the real place.
For every old local point at any $p\notin S_*$,
the normalized class
\[
             \mathcal A=\bigl((T-c)/(R-c),2\bigr)
                                  \in\operatorname{Br}(C_0)
 \tag{A.90.15}
\]
has zero evaluation.
This follows directly from the three exhaustive
valuation cases in A.85: the character is unramified;
a pole of $T$ has order divisible by four; and
reduction to $T=c$ gives a root tuple whose field
contains $\sqrt2$.  The constant normalization
has zero evaluation outside $S_*$ by its inclusion
in the definition of that set.

\emph{The controlled covers and the refinement map.}
At a stage of the construction the cover has equations
\[
 C_j:\quad T-c_i=a_{i,j}X_i^{n_{i,j}},\qquad
          4\mid n_{i,j},\qquad \prod_iX_i\ne0,\quad T\ne c.
 \tag{A.90.16}
\]
We maintain a finite set $S_j\supset S_*$.
It contains all coefficient, puncture-difference,
denominator, and exponent primes for the current
cover, and all ramified primes of the fixed $B_i$.
At every $p\in S_j$ retain an actual point $P_{j,p}$
of (A.90.16), with all three coordinate-root
evaluations (A.90.14) zero.
At the real place retain $T=R$ and positive root
coordinates.
At every $p\notin S_j$ there is a local point of
(A.90.16) with all three root coordinates units.
Except for the two coordinates $X_1,X_2$ at the
initial dyadic point until their first refinements,
all retained root coordinates are units as well.

These assertions hold initially: enlarge $S_0$
as necessary and use the stated points of $C_0$.
At places outside $S_0$, $P_0$ has unit coordinates.
We now prove a step that refines any chosen
coordinate $i$ by any integer $h\ge2$.
Write $n_i=n_{i,j}$ and $a_i'=a_{i,j}$ for this
paragraph.  We will choose $u\in\Q_{>0}$ and use
the exact map
\[
 \begin{split}
 C_{j+1}&\longrightarrow C_j,\\
 T&\longmapsto T,\qquad
 X_i=uY_i^h,\qquad X_k=Y_k\ (k\ne i),\\
 n_{i,j+1}&=hn_i,\qquad
 a_{i,j+1}=a_i'u^{n_i},\\
 n_{k,j+1}&=n_{k,j},\qquad
 a_{k,j+1}=a_{k,j}\quad(k\ne i).
 \end{split}
 \tag{A.90.17}
\]
The equations verify this map by direct substitution.
In particular $a_{i,j+1}/a_{i,j}$ is an
$n_{i,j}$th power in $\Q^\times$.

All coefficient ratios $a_{i,j}/a_i$ are fourth
powers in $\Q^\times$ by induction.
The boundary radicands at $T=c_i$ are now
$(c_i-c_k)/a_{k,j}$, with exponent $n_{k,j}$ for
$k\ne i$.
Their squareclasses are the fixed squareclasses
in (A.90.12).
Equations (A.90.3)--(A.90.10) show that their
common abelian fibre field is exactly the same
field $B_i$ at every stage.
This identifies the actual boundary characters
used below, without a constancy assertion for
the root functions $X_i$.

\emph{A coefficient-independent bound for the new cover.}
For three mixed exponents $\boldsymbol n=(n_1,n_2,n_3)$
put $D=\prod_i n_i$ and $N=\operatorname{lcm}_i n_i$.
The geometric Kummer cover has degree $D$.
For example, the inertia loop at $c_i$ acts as a
generator in the $i$th root coordinate and trivially
in the other two, so these three cyclic inertia
groups generate their direct product.
The inertia order at infinity is $N$, since its
diagonal element has order $\operatorname{lcm}_i n_i$.
Riemann--Hurwitz and removal of the fibre at $c$
give the exact genus and geometric boundary count
\[
 \begin{split}
 2g-2&=-2D+\sum_{i=1}^3D(1-1/n_i)+D(1-1/N),\\
 b&=\sum_{i=1}^3D/n_i+D/N+D .
 \end{split}
 \tag{A.90.18}
\]
The last $D$ points are the unramified fibre at $c$.
In particular $2g\le2D+2$ and $b\le5D$.
The explicit bound
\[
                       B(\boldsymbol n)=(8D+8)^2
 \tag{A.90.19}
\]
therefore ensures $p+1-2g\sqrt p-b>0$ for
$p>B(\boldsymbol n)$.
At a prime where the coefficients and puncture
differences are units and $p\nmid N$, this is a
good tame reduction with the same geometric cover.
Hasse--Weil supplies an interior residue point,
and its nonzero coordinate derivatives give a
unit-coordinate local point by Hensel lifting.
The trace and weight results used for this bound
are the same precise inputs as in A.48 and A.83,
\cite[(1.5.1), Theorem (1.6)]{AgentDeligneWeilI}
and \cite[Proposition 14.2]{AgentMilneLEC}.

\emph{The exact local root targets and the ray class.}
Before choosing $u$, enlarge $S_j$ to a finite set
$S$ containing every prime at most the new bound
(A.90.19), every prime divisor of the new
exponents, every prime at most $j+1$, and the
ramified primes of the splitting field of the
boundary fibre at $T=c_i$.
Retain the previously chosen points at $S_j$.
At a new prime in $S\setminus S_j$ choose an old
unit-coordinate local point, which exists by
the stage invariant.  All the $B_k$ are
unramified there, so every symbol
$(X_k,\chi)$ is zero: the reciprocity image of a
unit acts trivially in an unramified extension.
Thus every newly retained point satisfies
(A.90.14) as well.
All these points avoid $T=c$ by (A.90.16).

Let $x_p=X_i(P_{j,p})$ for $p\in S$ and let
$x_\infty>0$ be the real root coordinate.
Set
\[
       d=\prod_{p\in S}p^{v_p(x_p)}>0,\qquad
       v_p^*=x_p/d\in\Z_p^\times .
 \tag{A.90.20}
\]
Choose integers $M_p>2v_p(h)+4$, enlarging them
to contain the conductors of all characters of
$B_i$ in
$M=\prod_{p\in S}p^{M_p}$.
Set $\mathcal R=\Q(\mu_M)\supset B_i$ and choose
the exact CRT class
\[
 \rho\in(\Z/M)^\times,\qquad
             \rho\equiv v_p^*\pmod {p^{M_p}}\quad(p\in S).
 \tag{A.90.21}
\]
For every $\chi\in\operatorname{Hom}
(\operatorname{Gal}(B_i/\Q),\Q/\Z)$,
the invariant gives
$\chi(\operatorname{rec}_p x_p)=0$ at each
$p\in S$ and also at the real place.
Global reciprocity for the principal idele $d$
has zero contributions outside $S$, since $d$
is a unit and $B_i$ is unramified there.
The quotient $x_\infty/d$ is positive.
Subtracting this principal-idele identity gives
\[
 \sum_{p\in S}\chi(\operatorname{rec}_p v_p^*)=0.
 \tag{A.90.22}
\]
With arithmetic Frobenius on unramified
uniformizers, the product of these local unit
reciprocity maps on $\mathcal R$ is $\rho^{-1}$,
as in A.83.
Hence (A.90.22) proves the exact condition
\[
                         \rho|_{B_i}=1.
 \tag{A.90.23}
\]
The argument uses the individual vanishing of
symbols in the old root $X_i$, established in
(A.90.14); it does not replace that root by its
even power $T-c_i$.

\emph{Producing the rational coefficient.}
The finite boundary algebra for $X_i=0$ on
$C_j$ is
\[
 \mathcal E_i=
 \Q[Z_k^{\pm1}:k\ne i]/
       (Z_k^{n_{k,j}}-(c_i-c_k)/a_{k,j}:k\ne i).
 \tag{A.90.24}
\]
Indeed, at this boundary the other coordinate
derivatives are nonzero and $X_i$ itself is a
local parameter; its residue fields are exactly
the field factors of (A.90.24).
Let $L_i$ be its splitting field,
$G_i=\operatorname{Gal}(L_i\mathcal R/\Q)$,
$\Omega_i$ its root-tuple set, and
$\pi:G_i\to\operatorname{Gal}(\mathcal R/\Q)$
the restriction map.
The actual stabilizers have
\[
 \mathcal R^{\langle\pi((G_i)_\omega):
                                      \omega\in\Omega_i\rangle}
  =\mathcal R\cap\bigcap_{\omega\in\Omega_i}\Q(\omega)
  =B_i .
 \tag{A.90.25}
\]
The last equality follows from (A.90.10), since
the intersection within $\mathcal R$ is abelian
and $B_i\subset\mathcal R$.
Together with (A.90.23) this expresses $\rho$ as
a finite product $\gamma_1\cdots\gamma_t$,
where each $\gamma_l$ is the image of an actual
root-tuple stabilizer.

Chebotarev supplies distinct primes $q_l$
outside $S$ with these Frobenius restrictions
and with a fixed root tuple on the boundary.
They can exceed any prescribed bound.
Choose one further distinct prime $q_*$ splitting
completely in $L_i\mathcal R$.
Then
\[
          u=d\,q_*\prod_{l=1}^tq_l>0,\qquad
          u/x_p\in1+p^{M_p}\Z_p\ (p\in S),
          \qquad v_{q_*}(u)=1.
 \tag{A.90.26}
\]
The same valuation-one property holds at each
of the other fresh primes.
The existence assertion is the finite-conjugacy-class
form of Chebotarev
\cite[Theorem 8.31]{AgentMilneANT}.
The construction uses a product of stabilizer
images; no single-prime realization of a
generated subgroup is assumed.

\emph{All local points after this exact refinement.}
For $p\in S$ keep $T$ and every other old root
coordinate exactly.
By (A.90.26) and Hensel's lemma there is a root
\[
                    Y_i=(x_p/u)^{1/h}\in1+p^3\Z_p .
 \tag{A.90.27}
\]
More precisely, the polynomial $Y^h-x_p/u$
at $Y=1$ has error valuation at least $M_p$
and derivative valuation $v_p(h)$;
the inequality $M_p>2v_p(h)+4$ gives a root
with $v_p(Y_i-1)\ge M_p-v_p(h)>3$.
Every element of $1+p^3\Z_p$ is a square in
$\Q_p$, including $p=2$.
Consequently every quadratic character of
$B_i$ vanishes on $Y_i$.
The other root coordinates and the fixed
fields $B_k$ are unchanged, so all three
symbol-zero invariants persist.
The new root is a unit.
At the real place use the positive $h$th root
and retain $T=R$.

At a fresh coefficient prime
$q\in\{q_*,q_1,\ldots,q_t\}$, the boundary
fixed point gives nonzero residue roots for
(A.90.24).
Take
\[
                       Y_i=1,\qquad
                       T=c_i+a_{i,j}u^{n_{i,j}} .
 \tag{A.90.28}
\]
All old coefficients and puncture differences
are units at $q$, and $v_q(u)=1$.
The equations for $Y_k$, $k\ne i$, reduce to
the nonzero boundary roots.
Since $q$ does not divide any exponent, the
derivatives are units and each residue root
lifts independently.
All three new root coordinates are units.
Also $T\ne c$ and $T\ne c_k$ for every $k$:
the differences $c_i-c$ and $c_i-c_k$ are units
whereas $T-c_i$ has positive valuation.
All fixed $B_k$ are unramified at $q$, so every
new coordinate-root symbol is zero.
Retain this point permanently and add $q$ to
the controlled set.

At every prime outside $S$ and these fresh
primes, all new coefficients, exponent factors,
and differences are units.
The prime exceeds the new bound (A.90.19).
That bound supplies a local point with unit
coordinates avoiding $T=c$.
This establishes the outside-set invariant
for the next stage.
Taking $S_{j+1}$ to be $S$ together with all
fresh primes proves every part of the
induction.
In particular the two initially nonunit roots
at two become units on their first respective
refinements and remain units thereafter.

\emph{Cofinality and the completed classes.}
Apply this step in cycles indexed by
$h=2,3,4,\ldots$.
In cycle $h$, refine coordinates $1,2,3$ in
that order, each by the same factor $h$.
After completing that cycle their common
exponent is
\[
                              n_1=n_2=n_3=4h! .
 \tag{A.90.29}
\]
Thus for every positive integer $m$ all later
exponents are divisible by $m$.
At each individual step (A.90.17) preserves the
old coefficient modulo its old exponent.
Consequently, for each $i$, the coefficient
classes determine a compatible element of
the full completion:
\[
       A_i\bmod m=[a_{i,j}]\in\Q^\times/\Q^{\times m}
          \quad\text{at any sufficiently late stage }j .
 \tag{A.90.30}
\]
This is well defined by the displayed power
identity in (A.90.17); (A.90.29) proves
cofinality for every modulus, not only for
powers of two.
At modulus four the value is $[a_i]$.

Every fixed finite prime belongs to $S_j$
eventually, by the explicit inclusion of
primes at most $j+1$.
After that time its value $T=t_p$ is retained
exactly by every map (A.90.17).
For every modulus $m$, at a sufficiently
late stage its local equations give
\[
       t_p-c_i=a_{i,j}X_{i,j,p}^{n_{i,j}}
             \quad\Longrightarrow\quad
       [t_p-c_i]=[a_{i,j}]
                              \text{ in }\Q_p^\times/\Q_p^{\times m}.
 \tag{A.90.31}
\]
Taking all moduli proves the full local
completion equality in (A.90.2).
It also proves that every local value is one
ordinary nonzero number: no convergence
assertion for a varying local coordinate is
needed.
The retained points at $S_*$ give $t_2=16/17$
and $t_p=R$ for $p\in S_*\setminus\{2\}$.
The positive real lifts give the stated
ordinary real coordinate and signs.

Each time coordinate $i$ is refined, the
extra prime $q_*$ in (A.90.26) has
\[
       v_{q_*}(a_{i,j+1})=n_{i,j}>0,\qquad
       v_{q_*}(a_{k,j+1})=0\quad(k\ne i).
 \tag{A.90.32}
\]
It is controlled immediately with all root
coordinates units.
At every subsequent refinement the rational
factor $u$ has valuation zero there, by
(A.90.20) and (A.90.26).
Thus these exact integer valuations persist
in the completed classes.
There are infinitely many distinct such
primes for each coordinate, proving the
infinite-support assertion and excluding
equality with any rational coordinate class.

\emph{The obstruction to the new completed coordinate.}
The compatible maps to $C_0$ show that the
local point with coordinate $t_p$ is an
actual point of the old quartic cover at
every prime.
At the primes in $S_*$ its normalized
evaluation (A.90.15) is zero except at two,
where it is $1/2$ by A.69.
At every prime outside $S_*$ it is zero by
the proved outside-set calculation following
(A.90.15), regardless of the local point
chosen in the induction.
At the real place the quaternion with
second entry two is always split.

Suppose a global completed class for
$t_p-c$ existed and represent its reduction
modulo two by $b\in\Q^\times$.
The quaternion algebra
\[
                         (b/(R-c),2)
                                   \in\operatorname{Br}(\Q)
 \tag{A.90.33}
\]
would have these same local invariants.
Their sum is $1/2$, contradicting global
Brauer reciprocity
\cite[Theorem 12.1.8]{AgentCTSkoro2019}.
This excludes the global squareclass and
therefore the completed class, proving the
theorem.
\end{proof}

The constructed tuple satisfies the finite local
conditions of A.85 with actual completed old
coordinates.
Its points outside the fixed set $S_*$ are chosen
through the successive covers and are not prescribed
to equal one rational base point.
The exact local equalities (A.90.31) and the
infinite valuations (A.90.32) specify this fact
without invoking an inverse limit of rational
points on the old curves.
The theorem proves failure of extension for this
one completed old tuple.
It does not assert that every completed tuple
fails to extend, or that no ordinary tuple exists
for all rational punctures simultaneously.

The finite certificate
\path{code/completed_old_obstruction_certificate.py}
checks mixed-exponent stabilizers and their
generated images, the exact boundary constants
and initial local symbols, the cofinal refinement
maps and stable valuation bookkeeping, the
genus bound, and the local root precisions.
Its output is retained beside the code.
It does not claim to execute the unbounded
Chebotarev prime construction.
